%%%%%%%%%%%%%%%%%%%%%%%%%%%%%%%%%%%%%%%%%%%%%%%%%%%%%%%%%%%%
%% Lecture 11
%%%%%%%%%%%%%%%%%%%%%%%%%%%%%%%%%%%%%%%%%%%%%%%%%%%%%%%%%%%%

\lecture
{11}
{Wednesday, 16 September 2026}
{Introduction to Probability}
{Dr. Nishant Chandgotia}

%============================================================
% Independence of Generated Sigma-Fields
%============================================================

\begin{theorem}{Independence of Generated $\sigma$-Fields}
{independent-to-independent}

Suppose that
\[
\mathcal{A}_1,\mathcal{A}_2,\ldots,\mathcal{A}_n
\]
are independent $\pi$-systems. Then
\[
\sigma(\mathcal{A}_1),\sigma(\mathcal{A}_2),
\ldots,\sigma(\mathcal{A}_n)
\]
are independent.

\end{theorem}

\begin{proof}

We first prove the following intermediate statement:

\[
(\ast)\qquad
\text{If }
\mathcal{A}_1,\ldots,\mathcal{A}_n
\text{ are independent, then }
\sigma(\mathcal{A}_1),\mathcal{A}_2,\ldots,\mathcal{A}_n
\text{ are independent.}
\]

Once $(\ast)$ is proved, we can apply the same argument successively
to $\mathcal{A}_2,\ldots,\mathcal{A}_n$.

\medskip

Fix
\[
A_i\in\mathcal{A}_i,
\qquad 2\leq i\leq n,
\]
and define
\[
F=A_2\cap A_3\cap\cdots\cap A_n.
\]

Define
\[
\mathcal{L}
=
\left\{
A:
\mathbb{P}(A\cap F)
=
\mathbb{P}(A)\mathbb{P}(F)
\right\}.
\]

We will show that $\mathcal{L}$ is a $\lambda$-system and that
\[
\mathcal{A}_1\subseteq\mathcal{L}.
\]

Then the $\pi$-$\lambda$ theorem will imply
\[
\sigma(\mathcal{A}_1)\subseteq\mathcal{L}.
\]

\medskip

\textbf{Step 1: $\Omega\in\mathcal{L}$.}

Since
\[
\Omega\cap F=F
\]
and
\[
\mathbb{P}(\Omega)=1,
\]
we have
\[
\mathbb{P}(\Omega\cap F)
=
\mathbb{P}(F)
=
\mathbb{P}(\Omega)\mathbb{P}(F).
\]

Hence
\[
\Omega\in\mathcal{L}.
\]

\medskip

\textbf{Step 2: Difference property.}

Suppose
\[
A,B\in\mathcal{L},
\qquad A\subseteq B.
\]

Then
\[
(B\setminus A)\cap F
=
(B\cap F)\setminus(A\cap F).
\]

Since
\[
A\cap F\subseteq B\cap F,
\]
the difference formula for a probability measure gives
\[
\mathbb{P}((B\setminus A)\cap F)
=
\mathbb{P}(B\cap F)
-
\mathbb{P}(A\cap F).
\]

Since $A,B\in\mathcal{L}$,
\[
\mathbb{P}(B\cap F)
=
\mathbb{P}(B)\mathbb{P}(F)
\]
and
\[
\mathbb{P}(A\cap F)
=
\mathbb{P}(A)\mathbb{P}(F).
\]

Therefore
\[
\begin{aligned}
\mathbb{P}((B\setminus A)\cap F)
&=
\mathbb{P}(B)\mathbb{P}(F)
-
\mathbb{P}(A)\mathbb{P}(F)\\
&=
\bigl(\mathbb{P}(B)-\mathbb{P}(A)\bigr)\mathbb{P}(F).
\end{aligned}
\]

Since $A\subseteq B$,
\[
\mathbb{P}(B\setminus A)
=
\mathbb{P}(B)-\mathbb{P}(A).
\]

Hence
\[
\mathbb{P}((B\setminus A)\cap F)
=
\mathbb{P}(B\setminus A)\mathbb{P}(F).
\]

Therefore
\[
B\setminus A\in\mathcal{L}.
\]

\medskip

\textbf{Step 3: Continuity from below.}

Suppose
\[
B_k\in\mathcal{L},
\qquad
B_k\uparrow B.
\]

Then
\[
B_k\cap F\uparrow B\cap F.
\]

Indeed,
\[
B_k\subseteq B_{k+1}
\quad\Longrightarrow\quad
B_k\cap F\subseteq B_{k+1}\cap F,
\]
and
\[
\begin{aligned}
\bigcup_{k=1}^{\infty}(B_k\cap F)
&=
\left(\bigcup_{k=1}^{\infty}B_k\right)\cap F\\
&=
B\cap F.
\end{aligned}
\]

By continuity from below,
\[
\mathbb{P}(B\cap F)
=
\lim_{k\to\infty}
\mathbb{P}(B_k\cap F).
\]

Since $B_k\in\mathcal{L}$,
\[
\mathbb{P}(B_k\cap F)
=
\mathbb{P}(B_k)\mathbb{P}(F).
\]

Therefore
\[
\begin{aligned}
\mathbb{P}(B\cap F)
&=
\lim_{k\to\infty}
\mathbb{P}(B_k)\mathbb{P}(F)\\
&=
\mathbb{P}(F)
\lim_{k\to\infty}\mathbb{P}(B_k).
\end{aligned}
\]

Again, by continuity from below,
\[
\lim_{k\to\infty}\mathbb{P}(B_k)
=
\mathbb{P}(B).
\]

Consequently,
\[
\mathbb{P}(B\cap F)
=
\mathbb{P}(B)\mathbb{P}(F).
\]

Hence
\[
B\in\mathcal{L}.
\]

Thus $\mathcal{L}$ is a $\lambda$-system.

\medskip

\textbf{Step 4: Show that $\mathcal{A}_1\subseteq\mathcal{L}$.}

Take
\[
A_1\in\mathcal{A}_1.
\]

Since
\[
A_i\in\mathcal{A}_i,
\qquad
1\leq i\leq n,
\]
and
\[
\mathcal{A}_1,\ldots,\mathcal{A}_n
\]
are independent, we have
\[
\mathbb{P}(A_1\cap A_2\cap\cdots\cap A_n)
=
\prod_{i=1}^{n}\mathbb{P}(A_i).
\]

Also, since
\[
\mathcal{A}_2,\ldots,\mathcal{A}_n
\]
are independent,
\[
\mathbb{P}(F)
=
\mathbb{P}(A_2\cap\cdots\cap A_n)
=
\prod_{i=2}^{n}\mathbb{P}(A_i).
\]

Therefore
\[
\begin{aligned}
\mathbb{P}(A_1\cap F)
&=
\mathbb{P}(A_1\cap A_2\cap\cdots\cap A_n)\\
&=
\prod_{i=1}^{n}\mathbb{P}(A_i)\\
&=
\mathbb{P}(A_1)
\prod_{i=2}^{n}\mathbb{P}(A_i)\\
&=
\mathbb{P}(A_1)\mathbb{P}(F).
\end{aligned}
\]

Thus
\[
A_1\in\mathcal{L}.
\]

Since $A_1$ was arbitrary,
\[
\mathcal{A}_1\subseteq\mathcal{L}.
\]

\medskip

\textbf{Step 5: Apply the $\pi$-$\lambda$ theorem.}

The collection $\mathcal{A}_1$ is a $\pi$-system and
$\mathcal{L}$ is a $\lambda$-system containing $\mathcal{A}_1$.

Therefore, by the $\pi$-$\lambda$ theorem,
\[
\sigma(\mathcal{A}_1)\subseteq\mathcal{L}.
\]

Hence, for every
\[
A_1\in\sigma(\mathcal{A}_1),
\]
we have
\[
\mathbb{P}(A_1\cap F)
=
\mathbb{P}(A_1)\mathbb{P}(F).
\]

Since
\[
F=A_2\cap\cdots\cap A_n,
\]
it follows that
\[
\begin{aligned}
\mathbb{P}(A_1\cap A_2\cap\cdots\cap A_n)
&=
\mathbb{P}(A_1)
\mathbb{P}(A_2\cap\cdots\cap A_n)\\
&=
\mathbb{P}(A_1)
\prod_{i=2}^{n}\mathbb{P}(A_i)\\
&=
\prod_{i=1}^{n}\mathbb{P}(A_i).
\end{aligned}
\]

Thus
\[
\sigma(\mathcal{A}_1),
\mathcal{A}_2,\ldots,\mathcal{A}_n
\]
are independent.

This proves $(\ast)$.

\medskip

\textbf{Step 6: Iterate the argument.}

We have shown that
\[
\sigma(\mathcal{A}_1),
\mathcal{A}_2,\ldots,\mathcal{A}_n
\]
are independent.

By permuting the order of the systems, we may apply $(\ast)$ again
and obtain
\[
\sigma(\mathcal{A}_1),
\sigma(\mathcal{A}_2),
\mathcal{A}_3,\ldots,\mathcal{A}_n
\]
as independent.

Repeating this argument successively gives
\[
\sigma(\mathcal{A}_1),
\sigma(\mathcal{A}_2),
\ldots,
\sigma(\mathcal{A}_n)
\]
as independent.

Therefore
\[
\boxed{
\sigma(\mathcal{A}_1),\ldots,\sigma(\mathcal{A}_n)
\text{ are independent}.
}
\]

\end{proof}


%============================================================
% Criterion for Independence of Random Variables
%============================================================

\begin{observation}{A $\pi$-System Criterion for Independence}
Let
\[
X_1,X_2:(\Omega,\mathcal{F},\mathbb{P})
\longrightarrow
(S,\mathcal{R})
\]
be random variables.

Suppose that $\mathcal{A}\subseteq\mathcal{R}$ is a $\pi$-system
such that
\(
\sigma(\mathcal{A})=\mathcal{R}.
\)

Then, to prove that $X_1$ and $X_2$ are independent, it is enough
to verify that
\[
\mathbb{P}(X_1\in A_1,\ X_2\in A_2)
=
\mathbb{P}(X_1\in A_1)
\mathbb{P}(X_2\in A_2)
\]
for all
\(
A_1,A_2\in\mathcal{A}.
\)

Indeed,
\[
\mathcal{A}
\]
generates the relevant $\sigma$-field, and the preceding
$\pi$-$\lambda$ argument extends the independence from
$\mathcal{A}$ to $\sigma(\mathcal{A})$.

In particular, when
\[
(S,\mathcal{R})
=
(\mathbb{R},\mathcal{B}(\mathbb{R})),
\]
we may take
\[
\mathcal{A}
=
\left\{
(-\infty,a]:
a\in\mathbb{R}
\right\},
\]
which is a $\pi$-system generating
$\mathcal{B}(\mathbb{R})$.

\end{observation}


%============================================================
% Grouping Independent Sigma-Fields
%============================================================

\begin{theorem}{Independence of Groups of Independent $\sigma$-Fields}
{thm-after-ind-ind}

Suppose that
\[
\mathcal{F}_{i,j},
\qquad
1\leq i\leq n,
\quad
1\leq j\leq m(i),
\]
are independent $\sigma$-fields, and let
\[
\mathcal{G}_i
=
\sigma\left(
\bigcup_{j=1}^{m(i)}
\mathcal{F}_{i,j}
\right).
\]

Then
\(
\mathcal{G}_1,\ldots,\mathcal{G}_n
\)
are independent.

\end{theorem}

\begin{proof}

For each $i$, define
\[
\mathcal{A}_i
=
\left\{
\bigcap_{j=1}^{m(i)}A_{i,j}
:
A_{i,j}\in\mathcal{F}_{i,j}
\right\}.
\]

Since each $\mathcal{F}_{i,j}$ is a $\sigma$-field, in particular
\[
\Omega\in\mathcal{F}_{i,j}.
\]

We first show that $\mathcal{A}_i$ is a $\pi$-system.

Take
\[
A_i=\bigcap_{j=1}^{m(i)}A_{i,j},
\qquad
B_i=\bigcap_{j=1}^{m(i)}B_{i,j},
\]
where
\[
A_{i,j},B_{i,j}\in\mathcal{F}_{i,j}.
\]

Then
\[
A_i\cap B_i
=
\bigcap_{j=1}^{m(i)}
(A_{i,j}\cap B_{i,j}).
\]

Since each $\mathcal{F}_{i,j}$ is a $\sigma$-field,
\[
A_{i,j}\cap B_{i,j}\in\mathcal{F}_{i,j}.
\]

Therefore
\[
A_i\cap B_i\in\mathcal{A}_i.
\]

Hence $\mathcal{A}_i$ is a $\pi$-system.

Moreover, for every
\[
A_{i,j}\in\mathcal{F}_{i,j},
\]
we have
\[
A_{i,j}
=
A_{i,j}
\cap
\bigcap_{k\neq j}\Omega.
\]

Since
\[
\Omega\in\mathcal{F}_{i,k},
\]
it follows that
\[
A_{i,j}\in\mathcal{A}_i.
\]

Thus
\[
\bigcup_{j=1}^{m(i)}\mathcal{F}_{i,j}
\subseteq
\mathcal{A}_i.
\]

On the other hand, every element of $\mathcal{A}_i$ is a finite
intersection of sets belonging to
\[
\bigcup_{j=1}^{m(i)}\mathcal{F}_{i,j}.
\]

Hence
\[
\mathcal{A}_i
\subseteq
\sigma\left(
\bigcup_{j=1}^{m(i)}\mathcal{F}_{i,j}
\right)
=
\mathcal{G}_i.
\]

Therefore
\[
\sigma(\mathcal{A}_i)=\mathcal{G}_i.
\]

\medskip

Now take
\[
A_i\in\mathcal{A}_i,
\]
where
\[
A_i
=
\bigcap_{j=1}^{m(i)}A_{i,j},
\qquad
A_{i,j}\in\mathcal{F}_{i,j}.
\]

Since all the $\mathcal{F}_{i,j}$ are independent,
\[
\begin{aligned}
\mathbb{P}\left(
\bigcap_{i=1}^{n}A_i
\right)
&=
\mathbb{P}\left(
\bigcap_{i=1}^{n}
\bigcap_{j=1}^{m(i)}
A_{i,j}
\right)\\
&=
\prod_{i=1}^{n}
\prod_{j=1}^{m(i)}
\mathbb{P}(A_{i,j}).
\end{aligned}
\]

For each $i$,
\[
A_i
=
\bigcap_{j=1}^{m(i)}A_{i,j},
\]
and hence, again by independence,
\[
\mathbb{P}(A_i)
=
\prod_{j=1}^{m(i)}
\mathbb{P}(A_{i,j}).
\]

Therefore
\[
\mathbb{P}\left(
\bigcap_{i=1}^{n}A_i
\right)
=
\prod_{i=1}^{n}\mathbb{P}(A_i).
\]

Thus
\[
\mathcal{A}_1,\ldots,\mathcal{A}_n
\]
are independent $\pi$-systems.

By Theorem~\ref{independent-to-independent},
\[
\sigma(\mathcal{A}_1),\ldots,\sigma(\mathcal{A}_n)
\]
are independent.

Since
\[
\sigma(\mathcal{A}_i)=\mathcal{G}_i,
\]
we conclude that
\[
\boxed{
\mathcal{G}_1,\ldots,\mathcal{G}_n
\text{ are independent}.
}
\]

\end{proof}


%============================================================
% Functions of Independent Random Variables
%============================================================

\begin{theorem}{Functions of Independent Random Variables}
{functions-of-independent-random-variables}

Suppose that, for
\(
1\leq i\leq n,
\qquad
1\leq j\leq m(i),
\)
the random variables
\(
X_{i,j}
\)
are independent.

Suppose also that
\(
f_i:\mathbb{R}^{m(i)}\to\mathbb{R}
\)
are measurable functions.

Then the random variables
\(
f_i(X_{i,1},\ldots,X_{i,m(i)}),
\qquad
1\leq i\leq n,
\)
are independent.

\end{theorem}

\begin{proof}

For each $i,j$, let
\[
\mathcal{F}_{i,j}
=
\sigma(X_{i,j})
\]
and define
\[
\mathcal{G}_i
=
\sigma\left(
\bigcup_{j=1}^{m(i)}
\mathcal{F}_{i,j}
\right).
\]

Since the random variables $X_{i,j}$ are independent, the
$\sigma$-fields
\[
\mathcal{F}_{i,j}
\]
are independent.

Hence, by Theorem~\ref{thm-after-ind-ind},
\[
\mathcal{G}_1,\ldots,\mathcal{G}_n
\]
are independent.

\medskip

Now fix $i$.

Since
\[
X_{i,j}
\]
is $\mathcal{F}_{i,j}$-measurable and
\[
\mathcal{F}_{i,j}\subseteq\mathcal{G}_i,
\]
each $X_{i,j}$ is $\mathcal{G}_i$-measurable.

Therefore
\[
(X_{i,1},\ldots,X_{i,m(i)})
\]
is $\mathcal{G}_i$-measurable.

Since $f_i$ is measurable, the composition
\[
Y_i
=
f_i(X_{i,1},\ldots,X_{i,m(i)})
\]
is $\mathcal{G}_i$-measurable.

Consequently,
\[
\sigma(Y_i)\subseteq\mathcal{G}_i.
\]

\medskip

It remains to prove that
\[
Y_1,\ldots,Y_n
\]
are independent.

Take arbitrary
\[
B_i\in\sigma(Y_i),
\qquad
1\leq i\leq n.
\]

Since
\[
\sigma(Y_i)\subseteq\mathcal{G}_i,
\]
we have
\[
B_i\in\mathcal{G}_i
\]
for every $i$.

But
\[
\mathcal{G}_1,\ldots,\mathcal{G}_n
\]
are independent. Hence
\[
\mathbb{P}\left(
\bigcap_{i=1}^{n}B_i
\right)
=
\prod_{i=1}^{n}\mathbb{P}(B_i).
\]

Since this holds for every choice
\[
B_i\in\sigma(Y_i),
\]
the $\sigma$-fields
\[
\sigma(Y_1),\ldots,\sigma(Y_n)
\]
are independent.

Therefore
\[
Y_1,\ldots,Y_n
\]
are independent.

Recalling that
\[
Y_i
=
f_i(X_{i,1},\ldots,X_{i,m(i)}),
\]
we obtain
\[
\boxed{
f_1(X_{1,1},\ldots,X_{1,m(1)}),
\ldots,
f_n(X_{n,1},\ldots,X_{n,m(n)})
\text{ are independent}.
}
\]

\end{proof}


%============================================================
% Product Distribution of Independent Random Variables
%============================================================

\begin{theorem}{Product Distribution of Independent Random Variables}
{thm-product-of-random-variable}

Suppose $X_1,\ldots,X_n$ are independent random variables and
$X_i$ has distribution $\mu_i$.

Then the random vector
\(
(X_1,\ldots,X_n)
\)
has distribution
\(
\mu_1\times\cdots\times\mu_n.
\)

\end{theorem}

\begin{proof}

Let
\[
X=(X_1,\ldots,X_n).
\]

The distribution of $X$ is the probability measure
$\mu$ on $\mathbb{R}^n$ defined by
\[
\mu(B)
=
\mathbb{P}(X\in B),
\qquad
B\in\mathcal{B}(\mathbb{R}^n).
\]

Consider the collection
\[
\mathcal{A}
=
\left\{
A_1\times\cdots\times A_n:
A_i\in\mathcal{B}(\mathbb{R})
\right\}.
\]

This is a $\pi$-system generating
$\mathcal{B}(\mathbb{R}^n)$.

For
\[
A=A_1\times\cdots\times A_n\in\mathcal{A},
\]
we have
\[
\begin{aligned}
\mu(A)
&=
\mathbb{P}
\left(
X_1\in A_1,\ldots,X_n\in A_n
\right).
\end{aligned}
\]

Since
\[
X_1,\ldots,X_n
\]
are independent,
\[
\begin{aligned}
\mu(A)
&=
\prod_{i=1}^{n}
\mathbb{P}(X_i\in A_i)\\
&=
\prod_{i=1}^{n}\mu_i(A_i).
\end{aligned}
\]

But the product measure satisfies
\[
(\mu_1\times\cdots\times\mu_n)
(A_1\times\cdots\times A_n)
=
\prod_{i=1}^{n}\mu_i(A_i).
\]

Therefore
\[
\mu(A)
=
(\mu_1\times\cdots\times\mu_n)(A)
\]
for every
\[
A\in\mathcal{A}.
\]

Since $\mathcal{A}$ is a $\pi$-system generating
$\mathcal{B}(\mathbb{R}^n)$, the uniqueness theorem for probability
measures implies
\[
\mu
=
\mu_1\times\cdots\times\mu_n.
\]

Hence
\[
\boxed{
\mathcal{L}(X_1,\ldots,X_n)
=
\mu_1\times\cdots\times\mu_n.
}
\]

\end{proof}


%============================================================
% Sum and Product of Random Variables
%============================================================

\section{Sum and Product of Random Variables}

%------------------------------------------------------------
% Sums of Independent Random Variables
%------------------------------------------------------------

\subsection{Sums of Independent Random Variables}

\begin{theorem}{Distribution Function of the Sum}
{distribution-function-sum}

Let $X$ and $Y$ be independent random variables with distribution
functions
\[
F(x)=\mathbb{P}(X\leq x),
\qquad
G(y)=\mathbb{P}(Y\leq y).
\]

Then
\[
\boxed{
\mathbb{P}(X+Y\leq z)
=
\int_{\mathbb{R}}F(z-y)\,dG(y).
}
\]

\end{theorem}

\begin{proof}

Define
\[
h(x,y)
=
\mathbf{1}_{\{x+y\leq z\}}.
\]

Let $\mu$ and $\nu$ denote the probability distributions of $X$
and $Y$, respectively. Thus
\[
\mu(A)=\mathbb{P}(X\in A),
\qquad
\nu(B)=\mathbb{P}(Y\in B).
\]

Since $X$ and $Y$ are independent, their joint distribution is
\[
\mu\times\nu.
\]

Therefore
\[
\begin{aligned}
\mathbb{P}(X+Y\leq z)
&=
\mathbb{P}
\left(
(X,Y)\in
\{(x,y):x+y\leq z\}
\right)\\
&=
\iint
\mathbf{1}_{\{x+y\leq z\}}
\,\mu(dx)\,\nu(dy).
\end{aligned}
\]

For fixed $y$,
\[
h(x,y)
=
\mathbf{1}_{(-\infty,z-y]}(x).
\]

Hence
\[
\begin{aligned}
\int h(x,y)\,\mu(dx)
&=
\int
\mathbf{1}_{(-\infty,z-y]}(x)\,\mu(dx)\\
&=
\mu((-\infty,z-y])\\
&=
\mathbb{P}(X\leq z-y)\\
&=
F(z-y).
\end{aligned}
\]

Since $h\geq0$, Tonelli's theorem allows us to integrate first
with respect to $x$. Thus
\[
\begin{aligned}
\mathbb{P}(X+Y\leq z)
&=
\iint
\mathbf{1}_{\{x+y\leq z\}}
\,\mu(dx)\,\nu(dy)\\
&=
\int
\left[
\int
\mathbf{1}_{\{x+y\leq z\}}
\,\mu(dx)
\right]\nu(dy)\\
&=
\int F(z-y)\,\nu(dy).
\end{aligned}
\]

Since $G$ is the distribution function of $Y$, integration with
respect to $\nu$ is conventionally written as integration with
respect to $dG(y)$. Therefore
\[
\boxed{
\mathbb{P}(X+Y\leq z)
=
\int_{\mathbb{R}}F(z-y)\,dG(y).
}
\]

\end{proof}


%------------------------------------------------------------
% Density of the Sum
%------------------------------------------------------------

\begin{theorem}{Density of the Sum}
{density-of-sum}

Suppose that $X$ has density $f$ and $Y$ has distribution function
$G$, and that $X$ and $Y$ are independent.

Then $X+Y$ has density
\[
h(x)
=
\int_{\mathbb{R}}f(x-y)\,dG(y).
\]

If $Y$ also has a density $g$, then
\[
h(x)
=
\int_{\mathbb{R}}f(x-y)g(y)\,dy.
\]

\end{theorem}

\begin{proof}

For any $z\in\mathbb{R}$, by the preceding theorem,
\[
\mathbb{P}(X+Y\leq z)
=
\int_{\mathbb{R}}F(z-y)\,dG(y).
\]

Since $X$ has density $f$,
\[
F(z-y)
=
\int_{-\infty}^{z-y}f(u)\,du.
\]

Making the change of variable
\[
x=u+y,
\]
so that
\[
u=x-y,
\qquad
du=dx,
\]
we obtain
\[
F(z-y)
=
\int_{-\infty}^{z}f(x-y)\,dx.
\]

Consequently,
\[
\mathbb{P}(X+Y\leq z)
=
\int_{\mathbb{R}}
\left[
\int_{-\infty}^{z}
f(x-y)\,dx
\right]dG(y).
\]

Since $f\geq0$, Tonelli's theorem allows us to interchange the
order of integration:
\[
\begin{aligned}
\mathbb{P}(X+Y\leq z)
&=
\int_{-\infty}^{z}
\left[
\int_{\mathbb{R}}
f(x-y)\,dG(y)
\right]dx.
\end{aligned}
\]

Define
\[
h(x)
=
\int_{\mathbb{R}}
f(x-y)\,dG(y).
\]

Then
\[
\mathbb{P}(X+Y\leq z)
=
\int_{-\infty}^{z}h(x)\,dx.
\]

Thus $h$ is a density of $X+Y$, and hence
\[
\boxed{
h(x)
=
\int_{\mathbb{R}}f(x-y)\,dG(y).
}
\]

Finally, if $Y$ has density $g$, then
\[
dG(y)=g(y)\,dy.
\]

Therefore
\[
\boxed{
h(x)
=
\int_{\mathbb{R}}f(x-y)g(y)\,dy.
}
\]

This is the convolution formula for the density of the sum of two
independent random variables.

\end{proof}


%============================================================
% End of Lecture
%============================================================

\lectureend
